\documentclass[12pt,a4paper]{article}

\usepackage[spanish,es-nodecimaldot]{babel}
\usepackage[utf8]{inputenc}
\usepackage[T1]{fontenc}
\usepackage{lmodern}
\usepackage{amsmath,amssymb}
\usepackage{graphicx}
\usepackage{booktabs}
\usepackage{hyperref}
\usepackage{enumitem}
\usepackage{geometry}
\usepackage{listings}
\usepackage{xcolor}

\geometry{margin=2.5cm}

\lstset{
  basicstyle=\footnotesize\ttfamily,
  breaklines=true,
  frame=single,
}

\title{Glosario esencial + preguntas de tribunal}
\date{}

\begin{document}
\maketitle

\section{Glosario mínimo para explicar el proyecto a principiantes}

\begin{description}
  \item[Flujo de red] Resumen estadístico de una comunicación (no payload completo).
  \item[Feature] Variable numérica usada por el modelo para decidir.
  \item[Etiqueta (\texttt{y})] Clase real (\texttt{0=BENIGN}, \texttt{1=ATTACK}).
  \item[Observación] Vector que recibe el agente (aquí 152 dimensiones).
  \item[Política] Regla aprendida para elegir acción.
  \item[Recompensa] Señal numérica que define qué decisiones son mejores.
  \item[FP/FN] FP: bloquear benigno por error. FN: permitir ataque por error.
  \item[Recall de ataque] Proporción de ataques detectados.
  \item[F1 de ataque] Equilibrio entre precisión y recall de ataque.
  \item[Domain shift] Diferencia entre distribución de entrenamiento y tráfico real.
  \item[Scaler] Normalización estadística de features.
  \item[RUN\_ID] Identificador único de una ejecución reproducible.
\end{description}

\section{Preguntas frecuentes del tribunal (con respuesta técnica breve)}

\subsection{``¿Por qué RL y no solo clasificación supervisada?''}

Porque el proyecto necesita optimizar costes asimétricos de seguridad de forma explícita en la
función de recompensa, no solo exactitud global.

\subsection{``¿Qué aporta el esquema canónico?''}

Permite un contrato de observación estable entre datasets y Phase~2 real. Sin eso, el agente vería
vectores incompatibles entre entrenamientos.

\subsection{``¿Qué significa realmente 152?''}

76 features canónicas + 76 valores de máscara de missingness.

\subsection{``¿Cómo evitas leakage?''}

Excluyendo IPs, timestamps, Flow IDs y puertos-proxy en el preprocesado; además Check~B comprueba
que el rendimiento colapsa al barajar etiquetas.

\subsection{``¿Tu mejor accuracy implica robustez real?''}

No necesariamente. El random split da un techo in-distribution; la validación por día/CSV muestra
dificultad de generalización más realista.

\subsection{``¿Está lista Phase 2 para producción?''}

No. Está lista para inferencia offline robusta y diagnóstico; bloqueo activo y robustez operativa
cerrada aún no.

\subsection{``¿Cuál es tu principal limitación actual?''}

La sensibilidad al domain shift en tráfico real.

\section{Estructura de explicación en 3 minutos (versión corta)}

\begin{enumerate}
  \item problema: decisión PERMIT/BLOCK con costes asimétricos
  \item solución: contrato canónico (76+76), entorno RL, QRDQN, validación rigurosa
  \item resultado: alto rendimiento offline + identificación honesta de límite real (domain shift en
        Phase~2)
\end{enumerate}

\section{Estructura de explicación en 10--12 minutos (versión defensa)}

\begin{enumerate}
  \item motivación y objetivos
  \item datasets y esquema canónico
  \item entorno RL, recompensa y entrenamiento
  \item validaciones A/B/C + leave-one-exact-CSV-out
  \item Phase~2 robusta y riesgos abiertos
  \item conclusiones, límites y próximos pasos
\end{enumerate}

\end{document}
